\documentclass[12pt]{beamer}
\usepackage{color}
\usepackage{xcolor}
%\usepackage{latexsym}
%\usepackage{unicode-math}
\usepackage{amsmath}
\usepackage[labelfont=bf]{caption}
\usepackage{graphicx}  %package graphic
%\usepackage{siunitx}
\usepackage{tikz}
%\usepackage{algorithmicx}
%\usepackage[noend]{algpseudocode}
\usetikzlibrary{automata, positioning, arrows}
\usetheme{Boadilla}   
\usepackage{xeCJK}
%\usepackage{array}
%\usepackage{tabularx}
%\usepackage{mathtools}
%\usepackage{listings}
%\usepackage{textcomp}
%\usepackage[T1]{fontenc}
%\usepackage{lmodern}
\usepackage{stmaryrd}
\usepackage{adjustbox}

%\setCJKmainfont{微軟正黑體} 

\setbeamerfont{title}{size=\Large,series=\bfseries}  % title size

\setbeamerfont{frametitle}{size=\large,series=\bfseries}  % frametitle size, also can size*=<pt>

% Item include picture
\setbeamertemplate{itemize item}   % First Level item
{\includegraphics[height=0.33cm]{../Figures/golden-earth-on-white}}

\setbeamertemplate{itemize subitem} % Second level item
{\includegraphics[height=0.31cm]{../Figures/golden-sun-on-white}}

\setbeamertemplate{itemize subsubitem} % Third Level item
{\includegraphics[height=0.27cm]{../Figures/golden-paw-on-white}}

\definecolor{darkgold}{rgb}{0.765 0.64 0.0} % for highlighted text in black-and-white slides

%\newtheorem{problem}{Problem}
\mode<presentation>{\newtheorem{algorithm}{Algorithm}}
\mode<article>{\newenvironment{algorithm}{}{}}
%\newtheorem{solution}{Solution}

%\renewcommand{\highlightb}{\highlightg}

%\newlength{\subtextwidth}
%\setlength{\subtextwidth}{11cm}

\newenvironment{cbox}{
    \begin{center}
    \begin{tabular}{|l|}
    \hline
    \begin{minipage}[t]{\subtextwidth}}
    {\vspace{.25ex}
    \end{minipage}
    \hline
    \end{tabular}
    \end{center} }


\mode<presentation>{\title{Antichains: Alternative Algorithms for LTL Satisfiability and Model-Checking}}
\mode<article>{\title{Algorithms 2019: Analysis of Algorithms}}
%\subtitle{(Based on [Manber 1989])}
\author{林宏陽}
%\institute[IM.NTU]{Department of Information Management\\ National Taiwan University}
%\date[Algorithms 2019]{\null}
\mode<presentation>{\date[SVVRL]{\null}}
\mode<article>{\date{\today}}

\begin{document}
\begin{frame}
	\maketitle
\end{frame}

\begin{frame}{Introduction}
	\begin{itemize}
		\item A model for an LTL formula over a set $P$ of propositions : an infinite word $\omega$ over the alphabet $\Sigma = 2^{P}$.
		\item An LTL formula $\phi$ defines a set of words $\llbracket \phi \rrbracket = \{\omega \in \Sigma^{\omega} \mid \omega \models \phi\}$.
		\item The satisfiability problem : if $\llbracket \phi \rrbracket$ is empty.
		\item The model-checking problem : given an omega-regular language $\mathcal{L}$, if $\mathcal{L} \subseteq \llbracket \phi \rrbracket$
		\item Given a LTL formula $\phi$, we can construct a NBW $\mathcal{A}_{\phi}$ such that $\text{L}_{\text{b}}(\mathcal{A}_{\phi}) = \llbracket \phi \rrbracket$
	\end{itemize}
\end{frame}

\begin{frame}{Introduction(cont'd)}
	\begin{itemize}
		\item Implemented in explicit way : worst-case exponential.
		\item Implemented in symbolic way : it scales better.
		\item Efficient algorithms to reason on large LTL formulas are highly desirable.
		\item A new efficient way to analyze sABW (the number of states and symbolic transitions are linear in the size of the LTL formula).
	\end{itemize}
\end{frame}

\begin{frame}{Linear Temporal Logic}
	\begin{itemize}
		\item $P$ is a finite set of propositions
		\item a Kripke structure over
$P$ is a tuple $\mathcal{K} = \langle Q, q_{\iota}, \rightarrow_{\mathcal{K}}, \mathcal{L} \rangle$
		\begin{itemize}
			\item $Q$ is a finite set of states	
			\item $q_{\iota} \in Q$ is the initial state
			\item $\rightarrow_{\mathcal{K}} \subseteq Q \times Q$ is a transition relation
			\item $L : Q \to 2^{P}$ is a labeling function
		\end{itemize}
		\item A run of $\mathcal{K}$ is an infinite sequence $\rho = q_{0}q_{1}\cdots$ such that $q_{0} = q_{\iota}$ and for all $i \geq 0$, $(q_{i}, q_{i+1}) \in \rightarrow_{\mathcal{K}}$
		\item Let $\mathcal{L}(\rho) = \mathcal{L}(q_{0})\mathcal{L}(q_{1})$ and define the language of K as $\text{L}(\mathcal{K}) =
\{\mathcal{L}(\rho) \mid \rho \text{ is a run of } \mathcal{K}\}$
	\end{itemize}
\end{frame}

\begin{frame}{Linear Temporal Logic(cont'd)}
	\begin{itemize}
		\item The LTL formulas over $P$ are defined by $\phi ::= p \mid \neg\phi \mid \phi \lor \phi \mid \mathcal{X}\phi \mid \phi \mathcal{U} \phi$ where $p \in P$.
		\item Given an infinite word $\omega = \sigma_{0}\sigma_{1}\cdots \in \Sigma^{\omega}$ where $\Sigma = 2^{P}$ , and an LTL formula $\phi$ over $P$, we say that $\omega$ satisfies $\phi$ (written $\omega \models \phi$) if and only if (recursively):
		\begin{itemize}
			\item $\phi \equiv p$ and $p \in \sigma_{0}$,
			\item or $\phi \equiv \neg\phi_{1}$ and $\omega \models \phi_{1}$,
			\item or $\phi \equiv \phi_{1} \lor \phi_{2}$ and $\omega \models \phi_{1}$ or $\omega \models \phi_{2}$,
			\item or $\phi \equiv \mathcal{X}\phi_{1}$ and $\sigma_{1}\sigma_{2}\cdots \models \phi_{1}$,
			\item or $\phi \equiv \phi_{1}\mathcal{U}\phi{2}$ and for some $k \in N$, $\sigma_{k}\sigma_{k+1}\cdots \models \phi_{2}$ and for all $i$, $0 \leq i < k$,
$\sigma_{i}\sigma_{i+1}\cdots \models \phi_{1}$.
		\end{itemize}
		\item $true$, $false$ and $\phi_{1} \land \phi_{2}$ can be derived from the definition.
		\item let $\Box\phi = true\mathcal{U}\phi$, $\Box\phi =
\neg\Diamond\neg\phi$, and $\phi_{1}\mathcal{R}\phi_{2} = \neg(\neg\phi_{1}\mathcal{U}\neg\phi_{2})$.
	\end{itemize}
\end{frame}

\begin{frame}{Linear Temporal Logic(cont'd)}
	\begin{itemize}
		\item The satisfiability-checking problem : whether $\llbracket \phi \rrbracket = \emptyset$.
		\begin{itemize}
			\item we say that $\mathcal{K}$ satisfies $\phi$ ($\mathcal{K} \models \phi$) if and only if $\text{L}(\mathcal{K}) \subseteq \llbracket \phi \rrbracket$
			\item that is for all runs $\rho$ of $\mathcal{K}$, we have $\mathcal{L}(\rho) \models \phi$
		\end{itemize}
		\item The model-checking problem : whether $\mathcal{K} \models \phi$
		\item Both problems are PSPACE-COMPLETE.
	\end{itemize}
\end{frame}

\begin{frame}{sABW}
	\begin{itemize}
		\item Given a finite set $Q$, let $\text{Lit}(Q) = Q \cup \{\neg q \mid q \in Q\}$ be the set of literals over $Q$.
		\item $\mathcal{B}^{+}(Q)$ be the set of positive boolean formulas over $Q$ (formulas built from elements in $Q\cup\{true, false\}$ using $\land$ and $\lor$)
		\item Given $R \subseteq Q$ and $\varphi \in \mathcal{B}^{+}(Q)$, $R \models \varphi$ if and only if we assign $true$ to the
elements of $R$ and $false$ to the elements of $Q \setminus R$ satisfies $\varphi$.
	\end{itemize}
\end{frame}

\begin{frame}{sABW(cont'd)}
	\begin{itemize}
		\item A symbolic alternating Büchi automaton over the set of propositions $P$ is a tuple $\mathcal{A} = \langle \text{Loc}, I, \Sigma, \delta, \alpha \rangle$ where
		\begin{itemize}
			\item Loc is a finite set of states (or locations)
			\item $I \in \mathcal{B}^{+}(\text{Loc})$.
				\begin{itemize}
					\item A set $s \subseteq \text{Loc}$ is initial if $s \models I$.
				\end{itemize}
			\item $\Sigma = 2^{P}$ is the alphabet.
			\item $\delta : \text{Loc} \to \mathcal{B}^{+}(\text{Lit}(P)\cup\text{Loc})$ is the transition function.
			\item $\alpha \subseteq \text{Loc}$ is the set of accepting states.
		\end{itemize}
	\end{itemize}
\end{frame}


\begin{frame}{sABW(cont'd)}
	\begin{itemize}
		\item A run of $\mathcal{A}$ on an infinite word $\omega = \sigma_{0}\sigma_{1}\cdots \in \Sigma^{\omega}$ is a DAG $T_{\omega} = \langle V, V_{\iota}, \to \rangle$ where:
		\begin{itemize}
			\item $V \subseteq \text{Loc} \times N$. A node $(\ell, \mathnormal{i})$ represents the location $\ell$ after the first $\mathnormal{i}$ letters of $\omega$ have been read by $\mathcal{A}$. Nodes of the form $(\ell, \mathnormal{i})$ with 	$\ell \in \alpha$ are called $\alpha$-nodes;
			\item $V_{\iota} \subseteq \text{Loc} \times \{0\}$ is such that $V_{\iota} \subseteq V$ and $\{\ell \mid (\ell, 0) \in V_{\iota}\} \models I$;
			\item $\to \subseteq V \times V$ is such that for all $(\ell, \mathnormal{i}) \in V$
				\begin{itemize}
					\item if $(\ell, \mathnormal{i}) \to (\ell', \mathnormal{i}')$ then $\mathnormal{i}' = \mathnormal{i} + 1$
					\item $\sigma_{i} \cup \{\ell' \mid (\ell, \mathnormal{i}) \to (\ell', \mathnormal{i} + 1)\} \models \delta(\ell)$
				\end{itemize}
		\end{itemize}
		\item A run $T_{\omega}$ on an infinite word $\omega$ is accepting if \textcolor{red}{all} its infinite paths visit $\alpha$-nodes infinitely often.
		\item An infinite word $\omega \in \Sigma^{\omega}$ is accepted by $\mathcal{A}$ if there \textcolor{red}{exists} an accepting run on it.
	\end{itemize}
\end{frame}

\newcommand{\ABWExample}[1]{
\tikzset{
->, % makes the edges directed
>=stealth, % makes the arrow heads bold
every edge/.style={draw, thick, black},
node distance=0.8cm and 2.3cm, % specifies the minimum distance between two nodes. Change if necessary.
every state/.style={thick, fill=gray!10}, % sets the properties for each ’state’ node
initial text=$ $, % sets the text that appears on the start arrow
}
	\begin{figure}[ht] % ’ht’ tells LaTeX to place the figure ’here’ or at the top of the page
	\centering % centers the figure
	\begin{adjustbox}{max width=#1\textwidth, max totalheight=#1\textheight}
		\begin{tikzpicture}
			\node[draw=none] (l) {};
			\node[state, accepting, above right = 0.6cm and 1.5cm of l] (l4) {$\ell_{4}$};
			\node[state, right = of l4] (l3) {$\ell_{3}$};
			\node[state, below right =0.6cm and 1.5cm of l] (l2) {$\ell_{2}$};
			\node[state, accepting, right = of l2] (l1) {$\ell_{1}$};
			\node[draw=none, below right =0.4cm and 1cm of l3] (ll) {};
			\draw (l) edge[bend right = 25] node{} (l4);
			\draw (l) edge[bend left = 25] node{} (l2);
			\draw (l4) edge[loop above] node{$p$} (l4);
			\draw (l3) edge[loop above] node{true} (l3);
			\draw (l2) edge[loop below] node{true} (l2);
			\draw (l1) edge[loop below] node{$\neg r$} (l1);
			\draw (l4) edge node[above]{true} (l3);
			\draw (l4) edge[loop, out=0, in=-50, looseness=6] (l4);
			\draw (l2) edge node[above]{$\neg p \land \neg r$} (l1);
			\draw (l3) edge node[above]{$p$} (ll);
		\end{tikzpicture}
	\end{adjustbox}
		\caption{Alternating automaton for $\varphi \equiv \neg(\Box\Diamond p \to \Box(\neg p \to \Diamond r))$}
	\end{figure}
}


\begin{frame}{example}
	\ABWExample{0.45}
	\begin{itemize}
		\item $\ell_{4}$ and $\ell_{3}$ check that $\phi_{1} \equiv \Box\Diamond p$ holds.
		\item $\ell_{2}$ and $\ell_{1}$ check that $\phi_{2} \equiv \Diamond(\neg p \land \neg r)$ holds.
		\item $\delta(\ell_{4}) = (p \land \ell_{4}) \lor (\text{true} \land \ell_{3} \land \ell_{4})$
		\item $\delta(\ell_{3}) = (\text{true}\land \ell_{3})\lor (p \land \text{true})$
	\end{itemize}
\end{frame}

\begin{frame}{sNBW}
	\begin{itemize}
		\item A nondeterministic Büchi automaton is an sABW $\mathcal{A} = \langle \text{Loc}, I, \Sigma, \delta, \alpha \rangle$ such that
		\begin{itemize}
			\item $I$ is a disjunction of locations.
			\item for all $\ell \in \text{Loc}$, $\delta(\ell)$ is a disjunction of $\varphi \land \ell'$ where $\varphi \in \mathcal{B}^{+}(\text{Lit}(P))$ and $ \ell' \in \text{Loc}$.
		\end{itemize}
		\item Define the reverse automaton of $\mathcal{A}$ as $\mathcal{A}^{-1} = \langle \text{Loc}, I, \Sigma, \delta^{-1}, \alpha \rangle$ where $\delta^{-1} = \{(\ell, \sigma, \ell') \mid (\ell', \sigma, \ell) \in \delta\}$.
	\end{itemize}
\end{frame}

\begin{frame}{Miyano-Hayashi construction}
	\begin{itemize}
		\item Transform an sABW into a sNBW.
		\item The state of sNBW is the form $\langle s, o \rangle$.
		\item $s$ is a set of states that maintains a whole level of a guessed run DAG.
		\item $o$ is a set of states that ``owe'' a visit to an accepting state.
		\item  The Büchi condition asks that $o$ gets empty infinitely often in order to ensure that \textcolor{red}{every} path of the run DAG visits accepting states infinitely often.
	\end{itemize}
\end{frame}

\begin{frame}{Miyano-Hayashi construction (cont'd)}
	
		Given an sABW $\mathcal{A} = \langle \text{Loc}, I, \Sigma, \delta, \alpha \rangle$ over $P$, let MH$(\mathcal{A}) = \langle Q, I^{\text{MH}}, \Sigma, \delta^{\text{MH}}, \alpha^{\text{MH}} \rangle$ be a sNBW where:
		\begin{itemize}
			\item $Q = 2^{\text{Loc}} \times 2^{\text{Loc}}$
			\item $I^{\text{MH}}$ is the disjunction of all the pairs $s$, $\emptyset$ such that $s \models I$
		
		\item $\delta^{\text{MH}}$ is defined for all $\langle s, o \rangle \in Q$:
		\begin{itemize}
			\item if $o \neq \emptyset$, then $\delta^{\text{MH}}(\langle s, o \rangle)$ is the disjunction of all the formulas $\varphi \land \langle s', o' \setminus \alpha \rangle$ with $\varphi \in \mathcal{B}^{+}(\text{Lit}(P))$ such that
			\begin{itemize}
				\item $o' \subseteq s'$
				\item $\forall \ell \in s \cdot \forall \sigma \subseteq P$ : if $\sigma \models \varphi$ then $\sigma \cup s' \models \delta(\ell)$
				\item $\forall \ell \in o \cdot \forall \sigma \subseteq P$ : if $\sigma \models \varphi$ then $\sigma \cup o' \models \delta(\ell)$
			\end{itemize}
			
			\item if $o = \emptyset$, then $\delta^{\text{MH}}(\langle s, o \rangle)$ is the disjunction of all the formulas $\varphi \land \langle s', o' \setminus \alpha \rangle$ with $\varphi \in \mathcal{B}^{+}(\text{Lit}(P))$ such that
			\begin{itemize}
				\item $\forall \ell \in s \cdot \forall \sigma \subseteq P$ : if $\sigma \models \varphi$ then $\sigma \cup s' \models \delta(\ell)$
			\end{itemize}
		\end{itemize}
		\item $\alpha^{\text{MH}} = 2^{\text{Loc}} \times \{\emptyset\}$
		\end{itemize}
	\begin{theorem}
		For all sABW $\mathcal{A}$, we have $\text{L}_{\text{b}}(\text{MH}(\mathcal{A})) = \text{L}_{\text{b}}(\mathcal{A})$.
	\end{theorem}
\end{frame}

\begin{frame}{Recall ABW}
	\ABWExample{1}
\end{frame}

\begin{frame}{Miyano-Hayashi construction example}
\tikzset{
->, % makes the edges directed
>=stealth, % makes the arrow heads bold
every edge/.style={draw, thick, black},
node distance=0.8cm and 2.3cm, % specifies the minimum distance between two nodes. Change if necessary.
every state/.style={thick, fill=gray!10}, % sets the properties for each ’state’ node
initial text=$ $, % sets the text that appears on the start arrow
}
	\begin{figure}[ht] % ’ht’ tells LaTeX to place the figure ’here’ or at the top of the page
	\centering % centers the figure
	\begin{adjustbox}{max width=0.8\textwidth, max totalheight=0.8\textheight}
		\begin{tikzpicture}
			\node[state, initial] (a) {$\langle \{\ell_{2}, \ell_{4}\}, \{\ell_{2}\}\rangle$};
			\node[state, above right = and 4.5cm of a] (b) {$\langle \{\ell_{2}, \ell_{3},\ell_{4}\}, \{\ell_{2}\}\rangle$};
			\node[state, accepting, below right = of a] (c) {$\langle \{\ell_{1}, \ell_{3},\ell_{4}\}, \emptyset\rangle$};
			\node[state, accepting, right = of c] (d) {$\langle \{\ell_{1}, \ell_{4}\}, \emptyset\rangle$};
			\node[state, right = of d] (e) {$\langle \{\ell_{1}, \ell_{3},\ell_{4}\}, \{\ell_{3}\}\rangle$};

			\draw (a) edge[loop above] node[left]{$\{p\}, \{p, r\}$} (a);
			\draw (a) edge[bend left = 15] node[above = 0.5cm]{$\emptyset, \{p\}, \{r\}, \{p, r\}$} (b);
			\draw (b) edge[bend left = 15] node[below = 0.3cm]{$\{p\}, \{p, r\}$} (a);
			\draw (b) edge[loop above] node{$\emptyset, \{p\}, \{r\}, \{p, r\}$} (b);
			\draw (b) edge node[right]{$\emptyset$} (c);
			\draw (a) edge node[below]{$\emptyset$} (c);
			\draw (c) edge[loop below] node{$\emptyset, \{p\}$} (c);
			\draw (c) edge node[below]{$\{p\}$} (d);
			\draw (d) edge[loop above] node{$\{p\}$} (d);
			\draw (d) edge[bend left = 15] node[above]{$\emptyset, \{p\}$} (e);
			\draw (e) edge[bend left = 15] node[below]{$\{p\}$} (d);
			\draw (e) edge[loop above] node{$\emptyset, \{p\}$} (e);
		\end{tikzpicture}
	\end{adjustbox}
		\caption{After the Miyano-Hayashi construction}
	\end{figure}
\end{frame}

\end{document}